\section{Data curation}
\label{sec:data_curation}

A benchmark is only as trustworthy as the data underneath it.
Literature-aggregated solubility databases inherit errors at every
stage --- unit mismatches between publications, silent cross-author
data copying, undetected extraction artifacts, subtle mis-labelling of
compounds --- and none of the downstream analyses (aleatoric floor,
tier construction, model evaluation) can be interpreted before these
have been audited out.  \scthree{} is constructed from
\textsc{BigSolDB}~v2.1~\citep{krasnov2025bigsoldb} ---
\numval{112\,465} measurements of mole-fraction solubility covering
\numval{1\,525} solute \SMILES, \numval{218} solvent names,
\numval{1\,687} literature sources, across
\numval{243.15}--\numval{425.77}~K --- by applying a three-phase
pipeline of \emph{raw audit}, \emph{canonicalization}, and a
\emph{cleaning waterfall}.  

THE REST I SUGGEST PUTTING IN SI:

This section walks through the
three phases of data curation in details and reports the final dataset on which the
aleatoric analysis (§\ref{sec:aleatoric}) and benchmark splits
(§\ref{sec:benchmark}) are built.

\subsection{Raw-data audit}
\label{sec:raw_audit}

Before any filter is applied, we run an audit on the raw data.  Schema is
clean in the key columns (zero \texttt{NaN} entries in
\texttt{SMILES\_Solute}, \texttt{SMILES\_Solvent}, temperature or mole
fraction); \numval{3\,187} rows (\numval{2.8\,\%}) have \texttt{NaN}-valued
\logS, which are referred to complex organic solvents, for which \textsc{BigSolDB} lacks reliable information about their density, hindering units conversion.  All 1\,525 solute \SMILES{} were parsed with \texttt{RDKit}; one solvent
\SMILES{} is invalid (the literal \texttt{"-"} placeholder used for polymer
rows). Measurement temperatures are heavily clustered at 10 standard values of the
form $x.15$~K from 278.15 to 323.15 (covering \numval{68.9\,\%} of all
rows); the \emph{per}-DOI row count is right-skewed with median
\numval{56}, so no single source dominates.

\paragraph{Internal consistency of {\logS}.}
We back-compute
\[
    \log_{10} S = \log_{10}\!\left(\frac{x}{1-x}\cdot\frac{\rho(T)\,\cdot\,10^{3}}{M_w}\right)
\]
for every non-\texttt{NaN} row where a linear solvent-density model
$\rho(T) = aT + b$ is available, and compare to \textsc{BigSolDB}'s reported
\logS.  Across \numval{109\,278} testable rows the residuals are essentially
floating-point noise: median $|\mathrm{residual}| = 9.9\times 10^{-6}$,
P99 = $1.3\times 10^{-3}$.  Only in \numval{8} rows the error exceed $0.01$ log~S, all of which belong to a single DOI (\verb|10.1016/j.molliq.2017.02.075|,
$\beta$-alanine in methanol, systematic $-0.36$ log~S offset).  This DOI
is added to our bad-DOI list (§\ref{sec:waterfall}).

\paragraph{Cross-DOI bit-exact copycat detection.}
Grouped by $(\texttt{Solute\_Canon},\texttt{Solvent\_Canon},\texttt{round}(T,1),\texttt{round}(\logS,4))$,
\numval{86} four-tuples are reported by two or more DOIs at identical
values to the fourth decimal.  This affects \numval{172} rows across
\numval{46} distinct DOIs; the worst offender pair
(\verb|10.1016/jct.2024.107365| $\leftrightarrow$ \verb|10.1016/j.molliq.2025.127552|)
shares \numval{29} bit-identical rows.  These serve as Stage~A of the
source-integrity analysis in §\ref{sec:source}; in such case no threshold needs to be
tuned because the evidence is deterministic.

\paragraph{Inconsistent compound naming.}
The raw \texttt{Compound\_Name} column contains inconsistencies that do
not affect data integrity because our pipeline keys on \SMILES.  One
notable case: \numval{261} rows share a single raw \SMILES
(\verb|Cc1c(Cl)cccc1Nc1ccccc1C(=O)O|), a single CAS (\numval{13710-19-5})
and a single PubChem CID (\numval{610479}) --- all correct for tolfenamic
acid --- but two contributing DOIs label them inconsistently as either
``Tolfenamic acid'' or ``Pentoxifylline'' (the latter is a different
molecule; the \SMILES{} and CAS prove it is tolfenamic acid).  We flag
the case but take no action.

\subsection{Canonicalization policy}
\label{sec:canonicalization}

Molecule identification in this dataset hinges on \SMILES{}
canonicalization.  We compare four candidate policies against the raw
data and measure their empirical impact before choosing one
(Table~\ref{tab:canonicalization}).

\begin{table}[h]
\centering\small
\caption{Candidate canonicalization policies for solute \SMILES{}.}
\label{tab:canonicalization}
\begin{tabular}{@{}clr@{}}
\toprule
Option & Definition & Unique solutes \\
\midrule
A & tautomer-enumerate + strip \texttt{@/@@} + strip \texttt{/\textbackslash} & 1\,506 \\
B & no tautomer         + strip \texttt{@/@@} + strip \texttt{/\textbackslash} & 1\,508 \\
C & no tautomer         + strip \texttt{@/@@}, keep \texttt{/\textbackslash} & 1\,510 \\
\textbf{D} & plain canonical, \texttt{isomericSmiles=True} (keep all stereo) & \textbf{1\,525} \\
\bottomrule
\end{tabular}
\end{table}

Options A--C all merge distinct compounds.  Two concrete failures of
stereo stripping: \emph{fumaric} acid (trans, \verb|O=C(O)/C=C/C(=O)O|)
and \emph{maleic} acid (cis, \verb|O=C(O)/C=C\C(=O)O|), which differ in
aqueous solubility by roughly $100\times$, become identical after
removing \texttt{/}/\texttt{\textbackslash}.  A concrete failure of
tautomer enumeration: a bicyclic norbornene anhydride is silently
\emph{aromatized} by the RDKit tautomer enumerator into a structurally
nonsensical ``aromatic diol''.

To choose between the remaining options we empirically audit every
Option-C merge group.  For each group, we find all
$(\texttt{solvent}, T)$ cells where $\ge 2$ raw-SMILES members report
measurements and compute $|\Delta \log_{10} S|$ at matched conditions.
If enantiomers genuinely had identical solubility in achiral solvents,
these deltas should lie at or below the aleatoric floor
($\sim 0.1$ log~S; §\ref{sec:aleatoric}).  We find the opposite:
\emph{10 of 15} Option-C merge groups disagree by a \emph{median} of
\numval{0.29}--\textbf{\numval{1.31}} log~S across 3--117 overlap pairs
(Table~\ref{tab:stereo_audit}), i.e.\ 5--20$\times$ the aleatoric floor.

\begin{table}[h]
  \centering\small
  \caption{Empirical audit of Option-C stereo-merge groups.
    Median / max $|\Delta \log_{10} S|$ across all $(\mathrm{solvent}, T)$
    cells where $\ge 2$ raw-SMILES members of the merged group have
    measurements.  The four ``silent'' merges (bottom) have no overlap
    cells and cannot be empirically tested; they are structurally
    analogous to the tested cases (all D/L pairs) and we expect similar
    disagreement.}
  \label{tab:stereo_audit}
\begin{tabular}{@{}lrrrr@{}}
\toprule
Merge group (common name)              & total rows & overlap pairs & median $|\Delta|$ & max $|\Delta|$ \\
\midrule
Brassinolide (2 stereoisomers)         & 146 &  36 & \textbf{1.31} & 2.01 \\
D / L-psicose                          & 179 &  36 & 0.61 & 0.89 \\
D / L-malic acid                       & 141 &  21 & 0.52 & 0.99 \\
D / L-tryptophan                       & 229 &  84 & 0.43 & 0.97 \\
D / L-norbornene anhydride             & 245 & 117 & 0.40 & 0.56 \\
L-leucine / mixed                      & 126 &   8 & 0.39 & 0.45 \\
D / L-tyrosine                         & 123 &  16 & 0.37 & 1.07 \\
N-acetyl-methionine D / L              & 188 &  45 & 0.37 & 0.73 \\
Ibuprofen (S vs mixed)                 & 106 &   1 & 0.29 & 0.29 \\
\midrule
L-isoleucine (silent)                  & 126 &   0 &   --- &   --- \\
Ofloxacin (silent)                     & 124 &   0 &   --- &   --- \\
Naproxen (silent)                      & 115 &   0 &   --- &   --- \\
Tartaric acid (silent)                 &  49 &   0 &   --- &   --- \\
\bottomrule
\end{tabular}
\end{table}

The main chemically interpretable explanation is that enantiomeric compounds can potentially be crystallized in non-uniform conditions, being thus different polymorphic modifications (for which differences in solubility are commonly observed).
The other potential causes may be (ii) lab~$\times$~chirality-label
correlation in measurement protocol, (iii) silent mislabelling of
racemate as L (or vice versa) at the \textsc{BigSolDB} extraction stage.
We cannot disambiguate these from the data, but \emph{any} of the three
implies that merging averages away a 0.4-log-unit signal we cannot
recover.  We therefore adopt Option~D: trust the raw \SMILES, keep
every stereoisomer distinct, and let chirality-blind featurizations pay
their cost in model error rather than hiding that cost inside the label.

\subsection{Manual corrections}
\label{sec:sergei}

Some errors in the source database are fixable rather than only
detectable.  The source-integrity analysis of §\ref{sec:source}
flagged \numval{22} DOIs with inter-lab deviation exceeding $0.6$
log~S against peer consensus; we reviewed these candidates manually against the original papers and,
in several cases, identified specific transcription bugs rather than
the source paper being an outlier.  Four classes of correction were applied pre-cleaning
(Table~\ref{tab:sergei}); together they restore \numval{58} rows of
genuine data that would otherwise be dropped as bad sources.

\begin{table}[h]
\centering\small
\caption{Manual corrections applied before the cleaning waterfall.}
\label{tab:sergei}
\begin{tabular}{@{}lll r@{}}
\toprule
Class & Correction & DOI & rows \\
\midrule
C1 & paracetamol / water: 4 specific $x$ values replaced with audited values
   & \verb|10.1016/j.molliq.2020.113867|  &  4 \\
C2 & $\log_{10} S$ shifted by $+1$ ($\times 10$ unit correction)
   & \verb|10.1016/j.fluid.2013.09.018|   & 20 \\
C2 & $\log_{10} S$ shifted by $+1$ ($\times 10$ unit correction)
   & \verb|10.1016/j.molliq.2013.06.011|  & 10 \\
C3 & ethanol $\leftrightarrow$ ethyl acetate labels swapped, \logS{} re-derived
   & \verb|10.1016/j.fluid.2015.07.038|   & 24 \\
\bottomrule
\end{tabular}
\end{table}

\subsection{Cleaning waterfall}
\label{sec:waterfall}

Canonicalization and manual corrections fix individual rows.  The
remaining cleaning is a seven-stage waterfall applied to the
canonicalized, corrected table (Table~\ref{tab:waterfall}).  Each
stage has a single, documentable purpose; the polymer and salt
filters are based on \SMILES{} validity and structure rather than a
hand-curated name list.

\begin{table}[h]
\centering\small
\caption{Cleaning waterfall.  Each stage is applied in order to the
manually corrected canonicalized data.}
\label{tab:waterfall}
\begin{tabular}{@{}llrr@{}}
\toprule
Step & Filter                                                    & Rows after & Removed \\
\midrule
    & Input (canonicalized, manually corrected raw)                & 112\,465  &      --- \\
W1  & remove \numval{9} bad DOIs (see text)                      & 112\,019  &      446 \\
W2  & invalid / polymer solvent \SMILES{} (\texttt{"-"})         & 111\,705  &      314 \\
W3  & salts / mixtures (\texttt{.} in solute \SMILES)             & 101\,663  &  10\,042 \\
W4  & $M_w > 1000$ Da                                            & 101\,634  &       29 \\
W5  & \logS{} recovered for 2\,617 \texttt{NaN} rows; \numval{14} unrecoverable
                                                                 & 101\,620  &       14 \\
W6  & $\log_{10} S \notin [-15, 2]$                              & 101\,575  &       45 \\
W7  & intra-DOI near-duplicate dedupe ($T$ rounded to 0.1~K)     & \textbf{101\,535}  & 40 \\
\bottomrule
\end{tabular}
\end{table}

The \textbf{nine bad DOIs} excluded in W1 are: five confirmed as
systematic outliers by the manual review of §\ref{sec:sergei}
(\verb|10.1016/j.fluid.2011.09.033|,
 \verb|10.1021/jced.9b00728|,
 \verb|10.1021/jced.4c00179|,
 \verb|10.1021/jced.6b00009|,
 \verb|10.1016/j.molliq.2022.119759|);
three identified by the Hall-of-Shame audit of §\ref{sec:source}
(\verb|10.1021/je4000718|,
 \verb|10.1016/s1004-9541(08)60201-3|,
 \verb|10.1016/molliq.2016.11.036|);
and the \logS{}-residual DOI of §\ref{sec:raw_audit}
(\verb|10.1016/molliq.2017.02.075|).

The \textbf{salt/mixture filter} drops \numval{10\,042} rows containing
a \verb|.| in the solute \SMILES{} --- \numval{193} distinct
multi-component forms.  These are pharmaceutically important but exhibit
fundamentally different dissolution behaviour from the neutral form
(different lattice energy, different solvation, sometimes different
dissolution mechanism); modelling salt and/or cocrystals solubility properly requires
salt-specific features and labels and is left to a future
salt-aware split.

The \textbf{\logS{} recovery} step (W5) uses a three-tier density
lookup: (i) exact $(\text{solvent}, T)$ match from \textsc{BigSolDB}'s
own density file; (ii) a linear $\rho(T) = aT+b$ fit if available;
(iii) a lookup in the \texttt{thermo} library for the long tail.  Of
\numval{2\,631} surviving \texttt{NaN} rows, \numval{2\,617} are
recovered this way; the remaining \numval{14} use solvents the library
cannot identify by name, CAS, or \SMILES{}.

The \textbf{intra-DOI dedupe} in W7 removes \numval{40} rows, all from
a single DOI (\verb|10.1016/j.jct.2020.106137|) that reports each
nominal temperature twice ($T = x.15$~K and $T = x.16$~K) with \logS{}
identical to the fourth decimal --- a data-entry rounding artifact.

THE FINAL CLEANED DATASET SHOULD BELONG TO MAIN TEXT PROBABLY.

\paragraph{Final cleaned dataset.}
After canonicalization (Option D), manual corrections, and the
seven-stage waterfall, \scthree{} contains:

\begin{center}\small
\begin{tabular}{@{}lr@{}}
\toprule
measurements                                                 & \textbf{101\,535} \\
unique solutes (canonical)                                   & \textbf{1\,327}   \\
unique solvents                                              & \textbf{206}      \\
unique DOIs                                                  & \textbf{1\,493}   \\
temperature range                                            & 243.15--425.77 K  \\
\logS{} range                                                 & $-7.56$--$+1.99$  \\
$(\mathrm{solute},\mathrm{solvent})$ pairs with $\geq 2$ DOIs & \textbf{746} (6.8\%) \\
\bottomrule
\end{tabular}
\end{center}
